\title{\bfseries Validação técnica da autoria de grafos de comportamento assistida por agentes de LLM}
\author{João Carlos Queiroz\\\small EducaOFF\\\small \href{mailto:contato@educaoff.com.br}{contato@educaoff.com.br}}
\date{}

\maketitle

\begin{abstract}
Este estudo avalia tecnicamente o processo de autoria de grafos de comportamento da EducaOFF, no qual três agentes de linguagem propõem caminhos corretos, erros com remediações e cadeias de dicas, enquanto o GraphForge monta a estrutura final. A análise principal utilizou 24 exercícios de frações na reta numérica, distribuídos em seis estados de quatro problemas e três réplicas, com o código implantado congelado e as referências CTAT isoladas até o fechamento das saídas. Dezessete dos 18 estados foram concluídos. O agente 3a apresentou recall concreto ordenado de 0 e correspondência exata da resposta final de 0,111 (IC95\% 0,028--0,222). O 3b recuperou 0,176 (0,113--0,242) dos valores errados únicos registrados na referência. No braço de capacidade do 3c, o sucesso estrito por problema foi 0,278 (0,167--0,403), a validade estrita condicional das cadeias foi 0,382 (0,257--0,507) e a taxa de vazamento literal da resposta foi 0,226 (0,172--0,285). O extrator reteve 22,1\% dos itens do 3a, 86,0\% dos erros do 3b e 82,9\% das dicas do 3c, mas eliminou \path{problemId} em todos os fluxos; a política operacional também omitiu o 3c nos 17 estados completos. O GraphForge reproduziu 34/34 pares de artefatos de forma idêntica. Um painel auxiliar de LLMs apresentou efeito teto e sensibilidade limitada, não oferecendo validação pedagógica independente. Os resultados sustentam uma validação técnica parcial: a montagem é reprodutível, mas há fragilidades de conteúdo, segurança e rastreabilidade que impedem recomendar uso autônomo. Eles não demonstram equivalência a especialistas nem eficácia educacional.

\smallskip\noindent\textbf{Palavras-chave:} sistemas tutores inteligentes; grafos de comportamento; CTAT; agentes de LLM; example-tracing; dicas; validação técnica.
\end{abstract}

\begin{otherlanguage}{english}
\renewcommand{\abstractname}{Abstract}
\begin{abstract}
This study technically evaluates EducaOFF's behavior-graph authoring process, in which three language-model agents propose correct paths, errors with remediation, and hint chains, while GraphForge assembles the final structure. The main analysis used 24 fraction number-line exercises arranged in six four-problem states and three replications, with the deployed code frozen and CTAT references isolated until outputs were closed. Seventeen of 18 states were completed. Agent 3a achieved an ordered concrete recall of 0 and an exact final-answer match of 0.111 (95\% CI 0.028--0.222). Agent 3b recovered 0.176 (0.113--0.242) of the unique wrong values recorded in the reference. In Agent 3c's capacity arm, strict success per problem was 0.278 (0.167--0.403), conditional strict validity among emitted chains was 0.382 (0.257--0.507), and literal answer leakage was 0.226 (0.172--0.285). The extractor retained 22.1\% of Agent 3a items, 86.0\% of Agent 3b errors, and 82.9\% of Agent 3c hints, but removed \path{problemId} from every flow; the operational policy also omitted Agent 3c in all 17 completed states. GraphForge reproduced 34/34 artifact pairs identically. An auxiliary LLM panel showed a ceiling effect and limited sensitivity, providing no independent pedagogical validation. The findings support partial technical validation: assembly is reproducible, but weaknesses in content, safety, and traceability preclude recommending autonomous use. They do not establish equivalence to experts or educational effectiveness.

\smallskip\noindent\textbf{Keywords:} intelligent tutoring systems; behavior graphs; CTAT; LLM agents; example tracing; hints; technical validation.
\end{abstract}
\end{otherlanguage}
